\chapter*{从 Transformer 到智能体}
\addcontentsline{toc}{chapter}{从 Transformer 到智能体}

在前面的章节中，我们从感知机一路走到了 Transformer，理解了大型语言模型的技术基础。
从本章开始，我们将进入一个全新的领域——如何“用好”大模型，让它真正为我们做事。
在进入具体内容之前，我们先梳理一下这项技术演进的时间脉络。

\section*{2017：Transformer 诞生}
Google 发表论文《Attention Is All You Need》，提出 Transformer 架构。
这个架构抛弃了 RNN 的序列处理方式，完全基于注意力机制实现并行计算，
为后来所有大语言模型奠定了基石。

\section*{2020：GPT-3 与少样本学习}
OpenAI 发布 GPT-3，拥有 1750 亿参数。GPT-3 最令人惊讶的能力是——
它不需要微调，仅凭提示中的几个示例（Few-shot）就能完成各种任务。
这让人们意识到，大模型本身就是一个通用的任务求解器，关键在于你怎么"提问"。

\section*{2022：ChatGPT 与提示词工程}
2022 年底，ChatGPT 横空出世，大语言模型从实验室走向大众。
与此同时，提示词工程（Prompt Engineering）成为一门显学——
人们发现，同一个问题换一种问法，模型的回答质量天差地别。
Zero-shot、Few-shot、Chain-of-Thought 等技术相继被提出和普及。

\section*{2023：思维链与 RAG}
思维链（CoT）技术进一步成熟，让模型能够处理复杂的推理任务。
与此同时，检索增强生成（RAG）技术兴起，解决了大模型知识截止和幻觉问题——
模型可以"翻书"回答问题，而不是仅凭记忆。
这一年，工具使用（Tool Use）的概念也被引入，模型开始能够调用外部 API。

\section*{2024--2025：智能体时代}
Agent（智能体）成为最热门的方向。模型不再只是回答问题，而是能够自主规划、
使用工具、执行代码、操作文件，完成复杂的实际任务。
Claude Codex、OpenAI Codex 等编码智能体相继出现，
AI 从"对话助手"进化为"行动执行者"。

\section*{技术演进的本质}

回顾这条时间线，我们可以看到一条清晰的演进逻辑：

\begin{center}
	\begin{tabular}{|c|c|c|}
		\hline
		\textbf{阶段} & \textbf{核心问题} & \textbf{解决方案}              \\
		\hline
		提示词工程       & 如何问对问题        & Zero-shot / Few-shot / CoT \\
		RAG         & 如何获取最新知识      & 检索 + 生成                    \\
		工具使用        & 如何与外部交互       & Function Calling           \\
		智能体         & 如何自主完成任务      & 规划 + 工具 + 记忆               \\
		\hline
	\end{tabular}
\end{center}

本质上，这是一个让大模型从"能说"到"能做"的过程。
提示词工程解决了"怎么问"的问题，RAG 解决了"知道什么"的问题，
工具使用解决了"能做什么"的问题，而智能体将这些能力整合在一起，
让 AI 真正成为一个能够独立完成任务的助手。

接下来，我们将逐一深入这些技术，理解它们的原理和实现。
